

Recently, the majority of machine learning paradigms are primarily driven by data for the acquisition of knowledge, fundamentally transforming decision-making processes across various fields. However, the indiscriminate use of data has intensified the conflict between data rights and user privacy. In response to these pressing issues, a range of pioneering regulatory frameworks has been proposed globally, including the European Union’s General Data Protection Regulation (GDPR)\cite{GDPR}, and the California Consumer Privacy Act (CCPA)\cite{CCPA1} in California.
% the Personal Information Protection and Electronic Documents Act (PIPEDA)\cite{PIPEDA}, and the Canadian
% Consumer Privacy Protection Act (CPPA)\cite{CCPA2}. 
Among these regulations, one of the most critical and contentious provisions is \emph{the right to be forgotten}\cite{kwak2017let}. To technically align the effectiveness of machine learning with individual data rights and privacy regulations, a revolutionary frontier—machine unlearning\cite{Cao2015Towards,Xu2023ACMMU} has emerged.
% 隐私保护的重要性，引出研究主题

Despite the significant advancements in traditional machine unlearning for independently distributed data, the growing complexity of machine learning systems has underscored the rising importance of graph-structured data across various domains, including social networks\cite{app10041327,ijcai2018p142}, recommendation systems\cite{su2024dcl,LightGCN, cai2023app_gnn_rec3}, biological networks\cite{NIPS2017_f5077839,10.1093/bioinformatics/bty294,10149537}, and more. Hence, GNNs\cite{2016Convolutional_ChebNet,velivckovic2017gat,hamilton2017graphsage} have emerged as powerful tools for learning from this interconnected data, empowering systems to capture relationships and dependencies within graphs.
% However, the widespread adoption of GNNs brings with it heightened concerns over data privacy and the ethical use of graph data. This complexity of graph-based data, characterized by the intricate connections between nodes and edges, introduces unique challenges for unlearning, as the removal of a single data element will inevitably affect its neighboring information.
Building on this foundation, GU has emerged as a critical extension of unlearning techniques, specifically designed to meet the needs of graph-based machine learning models. Recently, a diverse array of GU methods has surfaced, paving the way for more responsible and privacy-compliant graph-based learning. Though these GU strategies have made strides from various perspectives, there are still deficiencies in achieving unity: (1) \textbf{Datasets}: The utilization of different datasets (i.e.,domain and scale) and rigidity of processing data regarding splitting and inference settings make it difficult to analyse the results. 
(2) \textbf{Backbones}: The lack of harmonization among GNN backbones between different methods (i.e., GCN\cite{kipf2016gcn} and SGC\cite{wu2019sgc}) creates barriers for fair comparisons. (3) \textbf{Experiments}: Typically, experiments are conducted on a single downstream task under one type of unlearning request, which prevents cross experiment of downstream tasks under different unlearning requests. Besides, the varying configurations of unlearning requests even within the same level and differing evaluation metrics in existing methods lead to divergent interpretations in experimental results. These challenges highlight the urgent need for a standardized benchmark to establish a unified and comprehensive understanding of the GU landscape.

To address the aforementioned challenges, we propose OpenGU in this paper, to the best of our knowledge, the first comprehensive benchmark specifically designed for graph unlearning. OpenGU integrates 16 recently proposed SOTA algorithms and 20 multi-domain datasets, covering node-level and edge-level tasks. OpenGU implemtents these components with unified APIs to address three types of unlearning requests(i.e., node, feature, and edge). Based on this design, we conduct an in-depth investigation of these GU algorithms, offering valuable insights in terms of four dimensions: effectiveness, efficiency and robustness. For \textbf{effectiveness}, OpenGU provides a comprehensive evaluation in two aspects: model updating and inference protection. We fairly assess the forgetting abilities for unlearning entities and the reasoning capabilities for retained data in order to determine whether the model achieves an effective trade-off between forgetting and reasoning. For \textbf{efficiency}, OpenGU provides insights into scalability by evaluating how well GU methods can handle increasing amounts of data and more complex graph structures, showing their practical applicability in real-world scenarios. Furthermore, we analyze the computational resources required by various graph unlearning algorithms, focusing on time and space complexity from both theoretical and empirical perspectives. For \textbf{robustness} we delve into the algorithms' capacity to maintain performance stability and integrity amid different deletion intensities. 

\textbf{Our contributions}. We propose OpenGU, the inaugural benchmark specifically tailored for GU domain. Our contributions are outlined as follows.
\begin{itemize}[leftmargin=*,noitemsep,topsep=0pt,parsep=0pt]
    \item \textbf{Comprehensive benchmark.} OpenGU integrates 16 SOTA algorithms and 20 popular multi-domain datasets, establishing a thorough evaluation framework for fair comparison. Moreover, OpenGU expands and standardizes experimental task requirements, encompassing unlearning requests and downstream tasks, thereby facilitating a code-level implementation of 3×2 cross experiment configurations.
    \item \textbf{Analytical Insights}. Through extensive empirical experiments centered around algorithms' forgetting capability and reasoning capability, we conduct a thorough analysis from the perspectives of effectiveness, scalability, efficiency and robustness, envisioning our conclusions as the catalyst for GU field.
    \item \textbf{Open-sourced benchmark library.} OpenGU is designed as an open-source benchmark library, providing researchers and practitioners with accessible tools and resources for expanding the exploration of GU. Additionally, comprehensive documentation and user-friendly interfaces fosters collaboration and innovation within the GU community. Our code is available at \textbf{XXXXXX}
\end{itemize}

\captionsetup[figure]{font=LARGE}
\begin{figure*}[h]
  \centering
  \includegraphics[width=1\linewidth,height=0.5\linewidth]{Pic/OpenGU (2).pdf}
  \caption{Overview of GU methods.}
  \label{figure1}
\end{figure*}

